\chapter{Introducción}
\label{cap:introduccion}

\section{Propósito y forma del informe}
\label{sec:proposito}

Este informe es el entregable de la Etapa 2, Diseño e Implementación Inicial, del curso
INFO1197 Trabajo de Título. Según lo acordado con el profesor guía el 5 de octubre de 2026,
cada etapa se entrega como un informe propio, organizado por los criterios de su pauta, con el
detalle en anexos y la evidencia enlazada al repositorio. Los capítulos \ref{cap:retro} a
\ref{cap:documentacion} siguen, en orden y con su título, los cuatro criterios de la Etapa 2
(20, 10, 55 y 15\,\%). Los enlaces remiten a destinos comprobados o archivos incluidos en el paquete local;
el anexo \ref{anexo:procedencia} distingue publicación, revisión local y acciones de entrega.

El trabajo reimplementa y refactoriza en C++ un instrumento científico heredado, el programa
DW-NOMINATE escrito en Fortran \parencite{poole2005spatial}, por mantenibilidad, operabilidad y
rendimiento. Ninguna de esas ventajas se da por demostrada: aquí son criterios de diseño, y se
miden en las etapas 3 y 4. La sustitución de las búsquedas heredadas por optimizadores modernos
es una intervención deliberada, aislada en un segundo motor. El caso chileno ejercita el
instrumento en condiciones operativas reales, y los datos de VoteView
\parencite{boche2018voteview} aportan una comparación adicional de confianza.

\section{Problema, objetivos y requisitos}
\label{sec:recap}

El planteamiento completo, con las partes interesadas, los criterios SMART y la evaluación de
alternativas, está en el Informe de Etapa 1 (\href{run:../evidence/informe-etapa-1.pdf}{Informe de Etapa 1}).

\paragraph{Problema.} La plataforma QueVotan estima las posiciones ideológicas de los
diputados con un mapa independiente por periodo, y dos mapas independientes no comparten
origen, orientación ni escala \parencite{herron2004scale}: el desplazamiento entre periodos no
es observable. DW-NOMINATE \parencite{poole2005spatial} estima todos los periodos en un marco
común, pero su implementación canónica es un programa Fortran de 2004, monohilo y no pensado
para ser auditado. Falta una implementación abierta, mantenible y auditable, y una
reimplementación no vale mientras no se establezca con qué fidelidad reproduce a la
referencia; la coincidencia de las salidas finales no basta, porque un defecto de indexación
abrió una distancia grande en log-verosimilitud con la primera dimensión prácticamente igual
\parencite{nieves2026jcc}.

\paragraph{Objetivos.} El objetivo general es desarrollar una implementación verificada en
C++ del modelo DW-NOMINATE que permita estimar posiciones ideológicas comparables en el
tiempo para la Cámara de Diputados de Chile. Los específicos son los de la Etapa 1, con los
plazos ajustados de la sección \ref{sec:cambios-e1}; la sección \ref{sec:propuesta-chile}
propone reformular OE3.
\begin{itemize}
  \item \textbf{OE1.} Implementar en C++ dos motores DW-NOMINATE, uno fiel al programa de
  referencia de 2004 y otro con optimizador moderno, con paridad de entradas y salidas
  respecto de la referencia.
  \item \textbf{OE2.} Verificar la fidelidad de la implementación frente a la referencia
  canónica sobre paneles estadounidenses y chilenos, midiendo la sensibilidad de la segunda
  dimensión al modelo temporal, al marco de coordenadas y a los valores iniciales.
  \item \textbf{OE3.} Estimar el desplazamiento de los partidos en la primera dimensión
  durante el 55.º periodo legislativo (2018--2022) con inferencia por bootstrap paramétrico,
  como caso de validación del instrumento.
\end{itemize}

\paragraph{Requisitos.} Se derivaron del artefacto construido, según lo establecido con el
propietario el 3 de septiembre de 2026 (cuadro \ref{tab:requisitos}); RNF09 se agregó después
de la Etapa 1. Las secciones \ref{sec:atributos} y \ref{sec:trazabilidad} los trazan al
diseño.

\begin{table}[htbp]
\centering
\caption{Requisitos funcionales y no funcionales del Informe de Etapa 1, con RNF09.}
\label{tab:requisitos}
\footnotesize
\begin{tabular}{L{1.1cm}L{4.9cm}L{1.1cm}L{5.0cm}}
\toprule
ID & Requisito funcional & ID & Requisito no funcional \\
\midrule
RF01 & Ingesta y validación de entradas & RNF01 & Fidelidad, dentro de una tolerancia
declarada por panel \\
RF02 & Valores iniciales declarados y reproducibles & RNF02 & Determinismo con cualquier
número de hilos \\
RF03 & Estimación estática & RNF03 & Reproducibilidad desde un clon limpio \\
RF04 & Estimación dinámica & RNF04 & Rendimiento, medido a un hilo sin alterar resultados \\
RF05 & Motor fiel & RNF05 & Portabilidad, Windows y Linux \\
RF06 & Motor moderno & RNF06 & Trazabilidad numérica \\
RF07 & Exportación en el marco del optimizador & RNF07 & Publicación \\
RF08 & Verificación en tres niveles & RNF08 & Control de versiones \\
RF09 & Informe de comparación & RNF09 & Mantenibilidad y auditabilidad \\
RF10 & Bootstrap paramétrico & & \\
RF11 & Inferencia sobre medianas de partido & & \\
RF12 & Sensibilidad al diseño & & \\
\bottomrule
\end{tabular}
\end{table}

\section{Los cuatro referentes}
\label{sec:referentes}

El informe distingue cuatro variantes del instrumento (cuadro \ref{tab:referentes}); el
original de 2004 y wmay nunca se tratan como un único «Fortran».

\begin{table}[htbp]
\centering
\caption{Los cuatro referentes que el informe nombra.}
\label{tab:referentes}
\footnotesize
\begin{tabular}{L{3.9cm}L{5.6cm}L{2.8cm}}
\toprule
Variante & Qué es & Rol en este trabajo \\
\midrule
Original 2004 (\archivo{DW-NOMINATE-2004.FOR}) & distribución de Poole y Rosenthal
\parencite{poole2005spatial}; SVD de IMSL; salvaguarda GRID & referencia histórica \\
wmay (\archivo{DW-NOMINATE-wmay.f}) & adaptación \parencite{wmay_dwnominate} con
\texttt{DGESDD} de LAPACK, sin GRID, acoplada a R & reproducción y evaluador común \\
Fiel C++ & traducción de las rutinas, corregida y configurada para compararse con wmay &
instrumento reimplementado \\
Moderno C++ & el mismo modelo con búsquedas de NLopt \parencite{johnson_nlopt} &
intervención \\
\bottomrule
\end{tabular}
\end{table}
